\subsection{Multiple Assignment and Unpacking}

Python assignment does not have to bind only one name at a time. A single assignment statement can bind several names from several produced values.

The simplest form is:

\begin{verbatim}
a, b = 10, 20
\end{verbatim}

This should be read as:

\begin{quote}
Evaluate the right side first, then bind the left-side names to the resulting values.
\end{quote}

A first inspection:

\begin{verbatim}
a, b = 10, 20

print(f"a = {a}")
print(f"b = {b}")

print()
print(f"type(a): {type(a)}")
print(f"type(b): {type(b)}")
\end{verbatim}

Possible output:

\begin{verbatim}
a = 10
b = 20

type(a): <class 'int'>
type(b): <class 'int'>
\end{verbatim}

The statement binds \texttt{a} to the integer object \texttt{10} and \texttt{b} to the integer object \texttt{20}. It does not create an array of variables. It is still name binding, only with several names in one statement.

A more explicit version uses parentheses on the right side:

\begin{verbatim}
a, b = (10, 20)

print(f"a = {a}")
print(f"b = {b}")
\end{verbatim}

The right side is a tuple object. Python unpacks that tuple into the names on the left.

\begin{verbatim}
pair = (10, 20)

print(f"pair: {pair}")
print(f"type(pair): {type(pair)}")

print("selected tuple methods:")
print(f"{'count':<10} {'count' in dir(pair)}")
print(f"{'index':<10} {'index' in dir(pair)}")
print(f"{'__len__':<10} {'__len__' in dir(pair)}")

print()
print(f"length: {len(pair)}")
print(f"count of 10: {pair.count(10)}")
print(f"index of 20: {pair.index(20)}")
\end{verbatim}

Possible output:

\begin{verbatim}
pair: (10, 20)
type(pair): <class 'tuple'>
selected tuple methods:
count      True
index      True
__len__    True

length: 2
count of 10: 1
index of 20: 1
\end{verbatim}

A tuple is an ordered collection object. For now, we only need this much: it can hold several values in order, and Python can unpack those values into several names.

\subsubsection{Parallel Assignment: Swapping Values with \texttt{a, b = b, a}}

A classic use of multiple assignment is swapping two names.

In many languages, swapping requires a temporary variable:

\begin{verbatim}
a = "Jerry"
b = "George"

temp = a
a = b
b = temp

print(f"a = {a}")
print(f"b = {b}")
\end{verbatim}

Possible output:

\begin{verbatim}
a = George
b = Jerry
\end{verbatim}

Python can express the same operation directly:

\begin{verbatim}
a = "Jerry"
b = "George"

print("before swap")
print(f"a = {a}")
print(f"b = {b}")

a, b = b, a

print()
print("after swap")
print(f"a = {a}")
print(f"b = {b}")
\end{verbatim}

Possible output:

\begin{verbatim}
before swap
a = Jerry
b = George

after swap
a = George
b = Jerry
\end{verbatim}

The important rule is that the right side is evaluated before the left side is rebound.

So this:

\begin{verbatim}
a, b = b, a
\end{verbatim}

means:

\begin{enumerate}
    \item read the current object referenced by \texttt{b};
    \item read the current object referenced by \texttt{a};
    \item create the pair of right-side results;
    \item bind \texttt{a} to the first result;
    \item bind \texttt{b} to the second result.
\end{enumerate}

Conceptually:

\begin{verbatim}
before:
    a ---> "Jerry"
    b ---> "George"

right side produces:
    ("George", "Jerry")

after:
    a ---> "George"
    b ---> "Jerry"
\end{verbatim}

A version with object identity:

\begin{verbatim}
a = "Nobody expects the Spanish Inquisition!"
b = "No soup for you."

print("before swap")
print(f"a id: {id(a)}")
print(f"b id: {id(b)}")

old_a_id = id(a)
old_b_id = id(b)

a, b = b, a

print()
print("after swap")
print(f"a id: {id(a)}")
print(f"b id: {id(b)}")

print()
print(f"a now has old b object: {id(a) == old_b_id}")
print(f"b now has old a object: {id(b) == old_a_id}")
\end{verbatim}

Possible output:

\begin{verbatim}
before swap
a id: 2144421957392
b id: 2144421957776

after swap
a id: 2144421957776
b id: 2144421957392

a now has old b object: True
b now has old a object: True
\end{verbatim}

The objects were not edited. The names were rebound.

This remains the central assignment model:

\begin{quote}
Assignment changes bindings between names and objects. It does not pour values into named storage boxes.
\end{quote}

\subsubsection{Sequence Unpacking and Arity Requirements}

Unpacking means taking values from an ordered object and binding them to names.

A tuple can be unpacked:

\begin{verbatim}
person = ("Bart", 10, "Springfield")

name, age, city = person

print(f"name: {name}")
print(f"age:  {age}")
print(f"city: {city}")

print()
print(f"type(person): {type(person)}")
print(f"type(name):   {type(name)}")
print(f"type(age):    {type(age)}")
print(f"type(city):   {type(city)}")
\end{verbatim}

Possible output:

\begin{verbatim}
name: Bart
age:  10
city: Springfield

type(person): <class 'tuple'>
type(name):   <class 'str'>
type(age):    <class 'int'>
type(city):   <class 'str'>
\end{verbatim}

A list can also be unpacked:

\begin{verbatim}
scores = [42, 87, 99]

first_score, second_score, third_score = scores

print(f"first:  {first_score}")
print(f"second: {second_score}")
print(f"third:  {third_score}")

print()
print(f"type(scores): {type(scores)}")
print(f"type(first_score): {type(first_score)}")
\end{verbatim}

Possible output:

\begin{verbatim}
first:  42
second: 87
third:  99

type(scores): <class 'list'>
type(first_score): <class 'int'>
\end{verbatim}

A list is also an ordered collection object. Unlike a tuple, a list is mutable. That distinction becomes important later. For this subsection, the point is simpler: Python can unpack ordered objects into names.

A short list inspection:

\begin{verbatim}
scores = [42, 87, 99]

print(f"scores: {scores}")
print(f"type(scores): {type(scores)}")

print("selected list methods:")
print(f"{'append':<10} {'append' in dir(scores)}")
print(f"{'count':<10} {'count' in dir(scores)}")
print(f"{'index':<10} {'index' in dir(scores)}")
print(f"{'pop':<10} {'pop' in dir(scores)}")

print()
print(f"length: {len(scores)}")
print(f"index of 87: {scores.index(87)}")
\end{verbatim}

Possible output:

\begin{verbatim}
scores: [42, 87, 99]
type(scores): <class 'list'>
selected list methods:
append     True
count      True
index      True
pop        True

length: 3
index of 87: 1
\end{verbatim}

Unpacking has an arity requirement. Arity means the number of values must match the number of receiving names.

This works:

\begin{verbatim}
x, y, z = [10, 20, 30]

print(x)
print(y)
print(z)
\end{verbatim}

This fails because there are too many right-side values:

\begin{verbatim}
x, y = [10, 20, 30]
\end{verbatim}

A safe demonstration:

\begin{verbatim}
try:
    x, y = [10, 20, 30]
except ValueError as err:
    print("unpacking failed")
    print(f"error type: {type(err)}")
    print(f"message: {err}")
\end{verbatim}

Possible output:

\begin{verbatim}
unpacking failed
error type: <class 'ValueError'>
message: too many values to unpack (expected 2)
\end{verbatim}

This also fails because there are too few right-side values:

\begin{verbatim}
x, y, z = [10, 20]
\end{verbatim}

Safe demonstration:

\begin{verbatim}
try:
    x, y, z = [10, 20]
except ValueError as err:
    print("unpacking failed")
    print(f"error type: {type(err)}")
    print(f"message: {err}")
\end{verbatim}

Possible output:

\begin{verbatim}
unpacking failed
error type: <class 'ValueError'>
message: not enough values to unpack (expected 3, got 2)
\end{verbatim}

The rule is exact unless extended unpacking is used:

\begin{quote}
Without a starred target, the number of receiving names must equal the number of produced values.
\end{quote}

Strings are also ordered objects, so they can be unpacked character by character:

\begin{verbatim}
a, b, c = "ABC"

print(a)
print(b)
print(c)

print()
print(f"type(a): {type(a)}")
\end{verbatim}

Possible output:

\begin{verbatim}
A
B
C

type(a): <class 'str'>
\end{verbatim}

This is legal, but it should be used carefully. Unpacking a string means unpacking its characters, not its words.

\begin{verbatim}
try:
    word1, word2 = "No soup"
except ValueError as err:
    print("unpacking failed")
    print(f"message: {err}")
\end{verbatim}

Possible output:

\begin{verbatim}
unpacking failed
message: too many values to unpack (expected 2)
\end{verbatim}

The string \texttt{"No soup"} has seven characters, including the space. It does not have two word-values from Python's point of view.

\subsubsection{Extended Unpacking with the Star Target (\texttt{*rest})}

Extended unpacking allows one left-side target to collect several values.

The starred name receives a list.

\begin{verbatim}
first, *rest = [10, 20, 30, 40]

print(f"first: {first}")
print(f"rest:  {rest}")

print()
print(f"type(first): {type(first)}")
print(f"type(rest):  {type(rest)}")
\end{verbatim}

Possible output:

\begin{verbatim}
first: 10
rest:  [20, 30, 40]

type(first): <class 'int'>
type(rest):  <class 'list'>
\end{verbatim}

This is useful when the first value has a special role and the remaining values should be kept together.

Example:

\begin{verbatim}
main_quote, *other_quotes = [
    "Nobody expects the Spanish Inquisition!",
    "No soup for you.",
    "D'oh!",
    "It is merely a flesh wound."
]

print(f"main quote: {main_quote}")
print(f"other quotes: {other_quotes}")

print()
print(f"type(main_quote): {type(main_quote)}")
print(f"type(other_quotes): {type(other_quotes)}")
\end{verbatim}

Possible output:

\begin{verbatim}
main quote: Nobody expects the Spanish Inquisition!
other quotes: ['No soup for you.', "D'oh!", 'It is merely a flesh wound.']

type(main_quote): <class 'str'>
type(other_quotes): <class 'list'>
\end{verbatim}

The starred target can also appear in the middle:

\begin{verbatim}
first, *middle, last = [10, 20, 30, 40, 50]

print(f"first:  {first}")
print(f"middle: {middle}")
print(f"last:   {last}")

print()
print(f"type(middle): {type(middle)}")
\end{verbatim}

Possible output:

\begin{verbatim}
first:  10
middle: [20, 30, 40]
last:   50

type(middle): <class 'list'>
\end{verbatim}

Or at the beginning:

\begin{verbatim}
*start, last = [10, 20, 30, 40]

print(f"start: {start}")
print(f"last:  {last}")

print()
print(f"type(start): {type(start)}")
print(f"type(last):  {type(last)}")
\end{verbatim}

Possible output:

\begin{verbatim}
start: [10, 20, 30]
last:  40

type(start): <class 'list'>
type(last):  <class 'int'>
\end{verbatim}

Only one starred target is allowed in a single unpacking assignment. This is invalid:

\begin{verbatim}
*left, *right = [10, 20, 30, 40]
\end{verbatim}

Python rejects it because two starred targets would make the split ambiguous. There would be no single obvious rule for deciding how many values go into \texttt{left} and how many go into \texttt{right}.

A safe syntax demonstration:

\begin{verbatim}
source_code = "*left, *right = [10, 20, 30, 40]"

try:
    compile(source_code, "<student-example>", "exec")
    print("compiled successfully")
except SyntaxError as err:
    print("syntax error")
    print(f"message: {err.msg}")
\end{verbatim}

Possible output:

\begin{verbatim}
syntax error
message: multiple starred expressions in assignment
\end{verbatim}

The starred target may receive an empty list:

\begin{verbatim}
first, *rest = [42]

print(f"first: {first}")
print(f"rest:  {rest}")

print()
print(f"type(rest): {type(rest)}")
print(f"length of rest: {len(rest)}")
\end{verbatim}

Possible output:

\begin{verbatim}
first: 42
rest:  []

type(rest): <class 'list'>
length of rest: 0
\end{verbatim}

This means the star target does not require at least one value. It collects whatever is left after the non-starred targets have received their required values.

However, the non-starred targets still have to be satisfied:

\begin{verbatim}
try:
    first, *middle, last = [42]
except ValueError as err:
    print("unpacking failed")
    print(f"message: {err}")
\end{verbatim}

Possible output:

\begin{verbatim}
unpacking failed
message: not enough values to unpack (expected at least 2, got 1)
\end{verbatim}

The rule is:

\begin{quote}
The starred target absorbs the flexible middle part, but the ordinary targets still require their own values.
\end{quote}

For assignment reasoning, multiple assignment and unpacking do not change the basic model. They only make the binding operation wider.

\begin{verbatim}
first, *middle, last = [10, 20, 30, 40]

print(f"{'name':<10} {'value':<15} {'type'}")
print(f"{'first':<10} {str(first):<15} {type(first)}")
print(f"{'middle':<10} {str(middle):<15} {type(middle)}")
print(f"{'last':<10} {str(last):<15} {type(last)}")
\end{verbatim}

Possible output:

\begin{verbatim}
name       value           type
first      10              <class 'int'>
middle     [20, 30]        <class 'list'>
last       40              <class 'int'>
\end{verbatim}

The practical summary is:

\begin{quote}
Multiple assignment evaluates the right side first, then binds several names. Unpacking distributes ordered values into names. Without a star target, the number of values must match exactly. With one star target, that name receives a list containing the flexible remainder.
\end{quote}